\uuid{cJzC}
\titre{Strucutre de sous-espace vectoriel}
\auteur{hanani ; mzari}
\organisation{exo7}
\licence{CC BY-NC-SA 4.0}
\datecreate{2020-01-20}
\niveau{L1}
\module{Algèbre}
\chapitre{Espace vectoriel}
\sousChapitre{Définition, sous-espace}
\difficulte{1}
\theme{Espaces vectoriels, Strucutre de sous-espace vectoriel}
\exo7qcmid{lille-2.espaces-vectoriels.004}

\contenu{
\begin{question}
\qclassification{106.01}
\qtags{motcle=Strucutre de sous-espace vectoriel}
Soit $E=\{(x,y,z) \in \Rr^3 \, ; \;  x-y+z = x+y-3z=0\}$, muni des opérations usuelles. Quelles sont les assertions vraies ?
\begin{answers}
\good{$E$ est non vide.}
\bad{$E$ n'est pas stable par addition.}
\good{$E$ est un espace vectoriel.}
\good{$E=\{(x,2x,x)  \, ; \,  x\in \Rr \}$.}
\end{answers}
\begin{explanations}
$E$ est un sous-espace vectoriel de $\Rr^3$, puisque $(0,0,0)\in E$, $E$ est stable par addition et multiplication par un scalaire.\\
En résolvant le système  : $\left\{\begin{array}{rcc}
x-y+z&=&0\\
x+y-3z&=&0\end{array}\right.$, on obtient : $E=\{(x,2x,x)  \, ; \,  x\in \Rr \}$.
\end{explanations}
\end{question}
}
